\subsection{Polynomial Rings, Divisibility, and Factorization}

Polynomial rings over fields behave like the integers in the ways linear algebra needs: division with remainder, gcds, Bezout identities, irreducible factorization, and principal ideals.

\begin{definition} \label{def:polynomial-degree-monic-splitting} \hlt{Polynomials over a Field}\\
Let $\mathbb{F}$ be a field.
\begin{enumerate}[label=(\roman*)]
\setlength{\itemsep}{0pt}
\item $\mathbb{F}[x]$ denotes the ring of polynomials in one variable over $\mathbb{F}$.
\item If $0\neq p\in\mathbb{F}[x]$ and $p(x)=a_0+a_1x+\cdots+a_nx^n$, $a_n\neq0$, then $\operatorname{lc}(p)\defeq a_n$, $\deg p\defeq n$.\\
Also, $p$ is \hlt{monic} iff $\operatorname{lc}(p)=1$. Set $\deg0=-\infty$.
\item A nonzero $p\in\mathbb{F}[x]$ \hlt{splits over $\mathbb{F}$} iff $p=\alpha\prod\limits_{i=1}^{m}(x-\lambda_i)$ for some $m\geq0$, $\alpha\in\mathbb{F}^{\times}$, and $\lambda_i\in\mathbb{F}$.
\end{enumerate}
\end{definition}

\begin{theorem} \label{thm:polynomial-division-algorithm} \hlt{Polynomial Division Algorithm}\\
Let $f,g\in\mathbb{F}[x]$ with $g\neq0$. Then $\exists!\,(q,r)\in\mathbb{F}[x]^2$ such that  $f=qg+r$, $r=0$ or $\deg r<\deg g$.
\begin{proof}
Let $d=\deg g$ and $\ell=\operatorname{lc}(g)$. Existence follows by induction on $\deg f$. If $f=0$ or $\deg f<d$, take $(q,r)=(0,f)$. Otherwise, let $n=\deg f$, $a=\operatorname{lc}(f)$, and $f_1=f-a\ell^{-1}x^{n-d}g$. Then $\deg f_1<n$.\\
By induction, $f_1=q_1g+r$ with $r=0$ or $\deg r<d$, hence $f=(q_1+a\ell^{-1}x^{n-d})g+r$.\\
For uniqueness, if $qg+r=q'g+r'$ with admissible remainders, then $(q-q')g=r'-r$.\\
If $q\neq q'$, then $\deg((q-q')g)\geq d$, while $\deg(r'-r)<d$, contradiction. Hence $q=q'$ and then $r=r'$.\\
\hlt{Structural Commentary}: Division works because the leading coefficient of $g$ is invertible in the field.\\
\hlt{Key Insight}: Cancel leading terms until the remainder is too small to divide again.
\end{proof}
\end{theorem}

\begin{definition} \label{def:polynomial-divisibility-gcd-coprime} \hlt{Divisibility, GCDs, and Coprimality}\\
For $p,q\in\mathbb{F}[x]$, write $p\mid q$ if $q=hp$ for some $h\in\mathbb{F}[x]$.\\
A \hlt{greatest common divisor} of $f,g\in\mathbb{F}[x]$, not both zero, is the unique monic polynomial $\gcd(f,g)$ such that:
\begin{enumerate}[label=(\roman*)]
\setlength{\itemsep}{0pt}
\item $\gcd(f,g)\mid f$ and $\gcd(f,g)\mid g$;
\item every common divisor of $f$ and $g$ divides $\gcd(f,g)$.
\end{enumerate}
The polynomials $f,g$ are \hlt{relatively prime} if $\gcd(f,g)=1$.
\end{definition}

\begin{theorem} \label{thm:polynomial-gcd-bezout} \hlt{Polynomial GCD and Bezout Identity}\\
For $f,g\in\mathbb{F}[x]$, not both zero, $\exists$ a unique monic $d\in\mathbb{F}[x]$ s.t. $d\mid f$, $d\mid g$, $h\mid f$, $h\mid g\Longrightarrow h\mid d$\\
Moreover, $d=\gcd(f,g)$ satisfies $d=af+bg$ for some $a,b\in\mathbb{F}[x]$.\\
Hence $\gcd(f,g)=1$ iff $\exists\,a,b\in\mathbb{F}[x]$ with $af+bg=1$.
\begin{proof}
If one of $f,g$ is $0$, the claim is immediate from the monic associate of the other.\\
Assume $f,g\neq0$ and run the Euclidean algorithm: $r_{i-1}=q_ir_i+r_{i+1}$, where $r_{i+1}<\deg r_i$, starting with $(r_0,r_1)=(f,g)$ and ending with $r_N\neq0$, $r_{N+1}=0$. Then $r_N\mid r_{N-1}$, and upward substitution gives $r_N\mid f,g$.\\
Conversely, any common divisor of $f,g$ divides every $r_i$, hence divides $r_N$. Thus $d\defeq\operatorname{lc}(r_N)^{-1}r_N$ is a monic gcd; uniqueness follows since two monic gcds divide each other.\\
Back-substitution gives $d=af+bg$ for some $a,b\in\mathbb{F}[x]$. Finally, $d=1$ gives $af+bg=1$; conversely, if $af+bg=1$, every common divisor divides $1$, so $d=1$.\\
\hlt{Structural Commentary}: The Euclidean algorithm produces both the gcd and its Bezout certificate.\\
\hlt{Key Insight}: GCD data is ideal data in $\mathbb{F}[x]$.
\end{proof}
\end{theorem}

\begin{definition} \label{def:irreducible-polynomial} \hlt{Irreducible Polynomials}\\
A nonconstant polynomial $p\in\mathbb{F}[x]$ is \hlt{irreducible over $\mathbb{F}$} if every factorization $p=ab$ in $\mathbb{F}[x]$ has $a$ or $b$ constant.
\end{definition}

\begin{proposition} \label{prop:irreducible-polynomial-prime-property} \hlt{Irreducible Polynomials Are Prime in $\mathbb{F}[x]$}\\
Let $p\in\mathbb{F}[x]$ be nonconstant. Then $p$ is irreducible over $\mathbb{F}$ iff $\forall f,g\in\mathbb{F}[x]$, $p\mid fg$ $\Longrightarrow$ $p\mid f$ or $p\mid g$.
\begin{proof}
\begin{enumerate}[label=\roman*.]
\setlength{\itemsep}{0pt}
\item[$(\Rightarrow)$] Assume $p$ irreducible and $p\mid fg$. If $p\mid f$, done. Otherwise $\gcd(p,f)=1$, so Theorem \ref{thm:polynomial-gcd-bezout} gives $a,b\in\mathbb{F}[x]$ with $ap+bf=1$. Multiplying by $g$ gives $apg+bfg=g$; since $p\mid apg$ and $p\mid bfg$, we get $p\mid g$.
\item[$(\Leftarrow)$] Assume the divisibility property and write $p=ab$. Since $p\mid ab$, either $p\mid a$ or $p\mid b$. If $p\mid a$, then $a=ph$, so $p=phb$ and cancellation gives $1=hb$, hence $b\in\mathbb{F}^{\times}$. Similarly, $p\mid b$ forces $a\in\mathbb{F}^{\times}$. Thus $p$ is irreducible.
\end{enumerate}
\hlt{Structural Commentary}: Over $\mathbb{F}[x]$, irreducible and prime are the same notion.\\
\hlt{Key Insight}: Bezout turns coprimality into a divisibility test.
\end{proof}
\end{proposition}

\begin{theorem} \label{thm:unique-factorization-polynomial-ring-field} \hlt{Unique Factorization in $\mathbb{F}[x]$}\\
For every nonconstant $f\in\mathbb{F}[x]$, $\exists$ $c\in\mathbb{F}^{\times}$ and monic irreducibles $p_1,\ldots,p_m\in\mathbb{F}[x]$ such that $f=c\prod\limits_{i=1}^{m}p_i$.\\
If $f=d\prod\limits_{j=1}^{n}q_j$ is another such factorization, then $m=n$, $c=d$, and $p_i=q_{\sigma(i)}$ for some $\sigma\in S_m$.
\begin{proof}
\begin{enumerate}[label=(\roman*)]
\setlength{\itemsep}{0pt}
\item \hlt{Existence}: Induct on $\deg f$. If $f$ is irreducible, write $f=\operatorname{lc}(f)\bigl(\operatorname{lc}(f)^{-1}f\bigr)$.\\
If $f=uv$ with $0<\deg u,\deg v<\deg f$, apply induction to $u$ and $v$ and concatenate the factors.
\item \hlt{Uniqueness}: Suppose $c\prod_{i=1}^{m}p_i=d\prod_{j=1}^{n}q_j$ with all $p_i,q_j$ monic irreducible. Since $p_1\mid\prod_{j=1}^{n}q_j$, iterating Proposition \ref{prop:irreducible-polynomial-prime-property} gives $p_1\mid q_j$ for some $j$. As $p_1,q_j$ are monic irreducible, $p_1=q_j$.\\
Cancel and continue by induction on $m$; the final scalar equality follows from leading coefficients.
\end{enumerate}
\hlt{Structural Commentary}: Polynomial factorization over a field mirrors prime factorization in $\Z$.\\
\hlt{Key Insight}: The prime property of irreducibles drives uniqueness.
\end{proof}
\end{theorem}

\begin{definition} \label{def:unit-associate-prime-irreducible-domain} \hlt{Units, Associates, Prime and Irreducible Elements}\\
Let $R$ be an integral domain.
\begin{enumerate}[label=(\roman*)]
\setlength{\itemsep}{0pt}
\item A \hlt{unit} is an element $u\in R$ with $uv=1$ for some $v\in R$.
\item Nonzero elements $a,b\in R$ are \hlt{associates}, written $a\sim b$, if $a=ub$ for some unit $u$.
\item A nonzero nonunit $p\in R$ is \hlt{prime} if $p\mid ab$ implies $p\mid a$ or $p\mid b$.
\item A nonzero nonunit $r\in R$ is \hlt{irreducible} if $r=ab$ implies $a$ or $b$ is a unit.
\end{enumerate}
\end{definition}

\begin{proposition} \label{prop:principal-ideals-divisibility-associates} \hlt{Divisibility and Principal Ideals}\\
Let $R$ be an integral domain and let $r,s\in R$ be nonzero.
\begin{enumerate}[label=(\roman*)]
\setlength{\itemsep}{0pt}
\item $r\in R^\times \Longleftrightarrow \langle r\rangle=R$.
\item $r\sim s \Longleftrightarrow \langle r\rangle=\langle s\rangle$.
\item $r\mid s \Longleftrightarrow \langle s\rangle\subseteq\langle r\rangle$.
\end{enumerate}
\begin{proof}
\begin{enumerate}[label=(\roman*).]
\setlength{\itemsep}{0pt}
\item[$(i)$]($\Rightarrow$) If $rv=1$, then $1\in\langle r\rangle$, so $\langle r\rangle=R$.
\item[$(i)$]($\Leftarrow$) If $\langle r\rangle=R$, then $1=ar$ for some $a\in R$, hence $r\in R^\times$.
\item[$(ii)$]($\Rightarrow$) If $r=us$ with $u\in R^\times$, then $r\in\langle s\rangle$ and $s=u^{-1}r\in\langle r\rangle$, so $\langle r\rangle=\langle s\rangle$.
\item[$(ii)$]($\Leftarrow$) If $\langle r\rangle=\langle s\rangle$, then $r=as$ and $s=br$ for some $a,b\in R$.\\
Thus $r=abr$, so $ab=1$ by cancellation; hence $a\in R^\times$ and $r\sim s$.
\item[$(iii)$]($\Rightarrow$) If $s=ar$, then every element of $\langle s\rangle$ is of the form $xs=xar\in\langle r\rangle$.
\item[$(iii)$]($\Leftarrow$) If $\langle s\rangle\subseteq\langle r\rangle$, then $s\in\langle r\rangle$, so $s=ar$ for some $a\in R$; hence $r\mid s$.
\end{enumerate}
\hlt{Structural Commentary}: Divisibility reverses inclusion of principal ideals.\\
\hlt{Key Insight}: Principal ideals translate multiplicative divisibility into order by containment.
\end{proof}
\end{proposition}

\begin{definition} \label{def:principal-ideal-domain-and-ufd} \hlt{PIDs and UFDs}\\
Let $R$ be an integral domain.
\begin{enumerate}[label=(\roman*)]
\setlength{\itemsep}{0pt}
\item $R$ is a \hlt{principal ideal domain} if every ideal $\mathcal{I}\trianglelefteq R$ has $\mathcal{I}=\langle a\rangle$ for some $a\in R$.
\item $R$ is a \hlt{unique factorization domain} if every $0\neq a\in R\setminus R^\times$ factors into irreducibles, uniquely up to reordering and associates.
\end{enumerate}
\end{definition}

\begin{theorem} \label{thm:integers-and-polynomial-ring-are-pids} \hlt{$\Z$ and $\mathbb{F}[x]$ Are Principal Ideal Domains}
\begin{enumerate}[label=(\roman*)]
\setlength{\itemsep}{0pt}
\item If $0\neq\mathcal{I}\trianglelefteq\Z$, then $\mathcal{I}=\langle d\rangle$ for $d=\min(\mathcal{I}\cap\mathbb{N}_{>0})$.
\item If $0\neq\mathcal{I}\trianglelefteq\mathbb{F}[x]$, then $\mathcal{I}=\langle m\rangle$ for the unique monic $m\in\mathcal{I}$ of minimal degree.
\end{enumerate}
\begin{proof}
\begin{enumerate}[label=(\roman*)]
\setlength{\itemsep}{0pt}
\item If $\mathcal{I}=0$, then $\mathcal{I}=\langle0\rangle$. Otherwise choose $d=\min(\mathcal{I}\cap\mathbb{N}_{>0})$. For $a\in\mathcal{I}$, write $a=qd+r$ with $0\leq r<d$. Since $r=a-qd\in\mathcal{I}$, minimality gives $r=0$, so $\mathcal{I}=\langle d\rangle$.
\item If $\mathcal{I}=0$, then $\mathcal{I}=\langle0\rangle$. Otherwise choose monic $m\in\mathcal{I}$ of minimal degree. For $p\in\mathcal{I}$, divide $p=qm+r$ with $r=0$ or $\deg r<\deg m$. Since $r=p-qm\in\mathcal{I}$, minimality gives $r=0$, so $\mathcal{I}=\langle m\rangle$.\\
If $m'$ is another monic minimal generator, then $m\mid m'$ and $m'\mid m$, hence $m=m'$.
\end{enumerate}
\hlt{Structural Commentary}: Euclidean division turns arbitrary ideals into one-generator ideals.\\
\hlt{Key Insight}: Choose the smallest positive integer or smallest-degree monic polynomial and divide by it.
\end{proof}
\end{theorem}

\begin{corollary} \label{cor:polynomial-generated-ideal-gcd} \hlt{GCDs Generate Polynomial Ideals}\\
For $0\neq p_1,\ldots,p_n\in\mathbb{F}[x]$ and monic $m=\gcd(p_1,\ldots,p_n)$, $\langle p_1,\ldots,p_n\rangle=\langle m\rangle$.
\begin{proof}
By Theorem \ref{thm:integers-and-polynomial-ring-are-pids}, $\langle p_1,\ldots,p_n\rangle=\langle g\rangle$ for a unique monic $g$. Since $p_i\in\langle g\rangle$, $g\mid p_i$ for all $i$.\\
If $q\mid p_i$ for all $i$, then $\langle p_1,\ldots,p_n\rangle\subseteq\langle q\rangle$, so $g\in\langle q\rangle$ and $q\mid g$. Hence $g=m$, and the claim follows.\\
\hlt{Structural Commentary}: In $\mathbb{F}[x]$, gcds and finitely generated ideals are the same data.\\
\hlt{Key Insight}: The gcd is the monic generator of the ideal generated by the polynomials.
\end{proof}
\end{corollary}

\begin{example} \label{ex:two-variable-polynomial-ring-not-pid} \hlt{A Polynomial Ring That Is Not a PID}\\
In $\mathbb{F}[x,y]$, the ideal $\langle x,y\rangle$ is not principal. If $\langle x,y\rangle=\langle p\rangle$, then $p\mid x$ and $p\mid y$, so $p\in\mathbb{F}^{\times}$. Hence $\langle x,y\rangle=\mathbb{F}[x,y]$, contradicting $1\notin\langle x,y\rangle$ since every element of $\langle x,y\rangle$ has zero constant term.
\end{example}

\begin{proposition} \label{prop:pid-ascending-chain-condition} \hlt{PIDs Satisfy the Ascending Chain Condition}\\
If $R$ is a PID and $\mathcal{I}_1\subseteq\mathcal{I}_2\subseteq\cdots$ is a chain of ideals, then $\exists N$ such that
\[
\mathcal{I}_j=\mathcal{I}_N\qquad\forall j\geq N.
\]
\begin{proof}
Let $\mathcal{I}=\bigcup_{j\geq1}\mathcal{I}_j$. By Proposition \ref{prop:intersections-and-chain-unions-of-ideals}, $\mathcal{I}$ is an ideal. Since $R$ is a PID, $\mathcal{I}=\langle a\rangle$. Choose $N$ with $a\in\mathcal{I}_N$. Then $\mathcal{I}=\langle a\rangle\subseteq\mathcal{I}_N\subseteq\mathcal{I}$, so $\mathcal{I}_j=\mathcal{I}_N$ for all $j\geq N$.\\
\hlt{Structural Commentary}: The generator appears at a finite stage.\\
\hlt{Key Insight}: Principal generation converts infinite ideal growth into a finite containment.
\end{proof}
\end{proposition}

\begin{theorem} \label{thm:pid-prime-irreducible-bezout} \hlt{Prime, Irreducible, and Bezout in a PID}\\
Let $R$ be a principal ideal domain.
\begin{enumerate}[label=(\roman*)]
\setlength{\itemsep}{0pt}
\item For $0\neq r\in R\setminus R^\times$, $r$ is irreducible iff $\langle r\rangle$ is maximal among proper ideals.
\item Every irreducible $r\in R$ is prime.
\item For $a,b\in R$, $\gcd(a,b)=1$ iff $\exists x,y\in R$ such that $xa+yb=1$.
\end{enumerate}
\begin{proof}
\begin{enumerate}[label=(\roman*)]
\setlength{\itemsep}{0pt}
\item If $r$ is irreducible and $\langle r\rangle\subseteq\mathcal{I}\subseteq R$, write $\mathcal{I}=\langle a\rangle$. Then $r=xa$, so irreducibility gives $\mathcal{I}=\langle r\rangle$ or $\mathcal{I}=R$.\\
Conversely, if $r=ab$ with $a,b\notin R^\times$, then $\langle r\rangle\subsetneq\langle a\rangle\subsetneq R$.
\item Let $r$ be irreducible and $r\mid ab$. By (i), $\langle r,a\rangle=\langle r\rangle$ or $R$. In the first case $r\mid a$.\\
In the second, $1=xr+ya$, so $b=xrb+yab$ and $r\mid b$.
\item If $xa+yb=1$, every common divisor of $a,b$ is a unit. Conversely, let $\langle a,b\rangle=\langle d\rangle$.\\
If $\gcd(a,b)=1$, then $d\in R^\times$, so $\langle a,b\rangle=R$ and $1=xa+yb$ for some $x,y\in R$.
\end{enumerate}
\hlt{Structural Commentary}: In a PID, factorization statements become ideal-containment statements.\\
\hlt{Key Insight}: Maximal principal ideals are the bridge from irreducible to prime.
\end{proof}
\end{theorem}

\begin{theorem} \label{thm:pid-is-ufd} \hlt{Every PID Is a UFD}\\
If $R$ is a PID, then $R$ is a UFD.
\begin{proof}
Let $0\neq r\in R\setminus R^\times$. If factorization into nonunits never terminates, obtain $s_1,s_2,\ldots$ with $s_{j+1}\mid s_j$ properly. By Proposition \ref{prop:principal-ideals-divisibility-associates},
\[
\langle s_1\rangle\subsetneq\langle s_2\rangle\subsetneq\cdots,
\]
contradicting Proposition \ref{prop:pid-ascending-chain-condition}; hence irreducible factorizations exist.\\
For uniqueness, let $p_1\cdots p_m=q_1\cdots q_n$ with all $p_i,q_j$ irreducible. By Theorem \ref{thm:pid-prime-irreducible-bezout}, $p_1$ is prime, so $p_1\mid q_j$ for some $j$; irreducibility gives $p_1\sim q_j$. Cancel associates and iterate. Thus $m=n$ and the factors agree up to order and associates.\\
\hlt{Structural Commentary}: ACC gives existence; irreducible-is-prime gives uniqueness.\\
\hlt{Key Insight}: PIDs have enough ideal control to force unique factorization.
\end{proof}
\end{theorem}

\begin{definition} \label{def:roots-of-unity-primitive-cyclotomic} \hlt{Roots of Unity and Cyclotomic Polynomials}\\
For $n\geq1$:
\begin{enumerate}[label=(\roman*)]
\setlength{\itemsep}{0pt}
\item The complex $n$th roots of unity are $U_n\defeq\{z\in\C:z^n=1\}$.
\item A root $\omega\in U_n$ is \hlt{primitive} iff $\operatorname{ord}(\omega)=n$; set $\Omega_n\defeq\{\omega\in U_n:\operatorname{ord}(\omega)=n\}$.
\item The $n$th \hlt{cyclotomic polynomial} is $\Phi_n(x)\defeq\prod\limits_{\omega\in\Omega_n}(x-\omega)$.
\item The number of primitive $n$th roots is $\varphi(n)\defeq|\Omega_n|$.
\end{enumerate}
\end{definition}

\begin{proposition} \label{prop:roots-of-unity-cyclic-and-cyclotomic-factorization} \hlt{Roots of Unity Factorization}\\
Let $\zeta_n=e^{2\pi i/n}$. Then $U_n=\langle\zeta_n\rangle$; $|U_n|=n$; $\Omega_n=\{\zeta_n^k:1\leq k\leq n,\ \gcd(k,n)=1\}$; $|\Omega_n|=\varphi(n)$.\\
Moreover, $U_n=\coprod\limits_{d\mid n}\Omega_d$ and $x^n-1=\prod\limits_{d\mid n}\Phi_d(x)$.
\begin{proof}
The roots of $x^n-1$ are precisely $\zeta_n^k$ for $0\leq k<n$, hence $U_n=\langle\zeta_n\rangle$ and $|U_n|=n$.\\
In a cyclic group, $\operatorname{ord}(\zeta_n^k)=n/\gcd(k,n)$, so $\zeta_n^k$ is primitive iff $\gcd(k,n)=1$; thus $|\Omega_n|=\varphi(n)$.\\
Each $z\in U_n$ has unique order $d\mid n$, giving $U_n=\coprod_{d\mid n}\Omega_d$.\\
Therefore $x^n-1=\prod\limits_{z\in U_n}(x-z)=\prod_{d\mid n}\prod\limits_{\omega\in\Omega_d}(x-\omega)=\prod\limits_{d\mid n}\Phi_d(x)$
\hlt{Structural Commentary}: Orders partition roots of unity by divisors.\\
\hlt{Key Insight}: Cyclicity turns root counting into order counting.
\end{proof}
\end{proposition}

\begin{proposition} \label{prop:cyclotomic-polynomials-integer-coefficients} \hlt{Cyclotomic Polynomials Have Integer Coefficients}\\
For every $n\geq1$, $\Phi_n(x)\in\Z[x]$ is monic.
\begin{proof}
Induct on $n$. For $n=1$, $\Phi_1(x)=x-1$.\\
Assume $\Phi_d\in\Z[x]$ for every proper divisor $d\mid n$, and set $P_n(x)\defeq\prod\limits_{\substack{d\mid n\\ d<n}}\Phi_d(x)\in\Z[x]$.\\
Then $P_n$ is monic, and Proposition \ref{prop:roots-of-unity-cyclic-and-cyclotomic-factorization} gives $x^n-1=\Phi_n(x)P_n(x)$.\\
Divide $x^n-1$ by $P_n$ in $\Z[x]$: $x^n-1=hP_n+r$, where $h,r\in\Z[x]$ and $r=0$ or $\deg r<\deg P_n$.\\
Uniqueness of division over $\C[x]$ gives $h=\Phi_n$ and $r=0$. Hence $\Phi_n\in\Z[x]$, and $\Phi_n$ is monic by definition.\\
\hlt{Structural Commentary}: Proper divisors isolate the next cyclotomic factor.\\
\hlt{Key Insight}: Monic division keeps the quotient in $\Z[x]$.
\end{proof}
\end{proposition}

\begin{lemma} \label{lem:cyclotomic-size-estimate-for-wedderburn} \hlt{A Cyclotomic Size Estimate}\\
For $n>1$ and $z\in\Z_{\geq2}$, $|\Phi_n(z)|>z-1$.
\begin{proof}
For $\omega\in\Omega_n$, $\omega\neq1$ and $|\omega|=1$, hence $|z-\omega|>z-1$. Thus $|\Phi_n(z)|=\prod\limits_{\omega\in\Omega_n}|z-\omega|>(z-1)^{\varphi(n)}\geq z-1$.\\
\hlt{Structural Commentary}: Primitive factors all lie outside the closest point $1$.\\
\hlt{Key Insight}: Each primitive root contributes a factor larger than $z-1$.
\end{proof}
\end{lemma}
